\section{Literature Review}

\subsection{The Persistent P/B Disparity: Defining the Korea Discount}
The ``Korea Discount'' refers to the chronic undervaluation of South Korean equities relative to global peers, measured through depressed price-to-book (P/B) and price-to-earnings (P/E) ratios. Across 48 countries from 1990 to 2023, Korean firms exhibit persistently lower P/B and P/E multiples---a finding that survives controls for firm size, profitability, growth, GDP per capita, financial market development, and geopolitical risk across all sub-periods \citep{kim_korean_market_longterm}. At the firm level, the high prevalence of Korean firms with book-to-market ratios above one is more consistent with systematic mispricing than fundamental risk, and a governance factor---rather than geopolitical or purely fundamental factors---is the cross-sectional predictor most strongly associated with this pattern \citep{jung_korea_discount_book_to_market}.

The literature is divided, however, on which root causes are most likely. Direct tests of governance, shareholder payout, and growth potential hypotheses find that corporate governance has no statistically significant relationship with PBR, while variables capturing stagnant growth prospects---low R\&D expenditure, high tangible-asset intensity, and firm maturity---exhibit positive correlations with low PBR \citep{yang_korea_discount_growth_vs_governance}. Higher shareholder payout tendencies are also associated with \emph{lower} PBR, counter to the shareholder-return argument. Consistent with the growth view, for large manufacturing firms (which account for most of the KOSPI) dividend expansion has weak or even negative effects on PBR, while R\&D investment is the more significant driver \citep{lee_korea_discount_pbr}.

These findings directly challenge the governance-reform interpretation, and this is not fully resolved by the existing literature. The distinction between cross-sectional correlations and causal identification through natural experiments, reviewed below, partly explains the divergence: natural experiments consistently implicate governance as a direct value-destroyer even where cross-sectional regressions cannot detect it. The Korea Discount is likely overdetermined: weak growth fundamentals and governance failures reinforce one another, and reforms addressing only one dimension will underperform.

North Korean geopolitical risk is a commonly cited but empirically weak alternative. Cross-country studies that control for geopolitical risk explicitly find the Korea Discount survives intact, consistent with markets treating North Korean threats as a chronic, well-priced background condition. Section~4.3 formalizes this test; we reject the geopolitical explanation and focus on the governance and fundamental factors reviewed below.

\subsection{Root Causes: Concentrated Ownership, Chaebols, and Tunneling}
The structural foundation of the Korea Discount lies in Korea's Chaebol system---family-controlled conglomerates that dominate the Korean economy. Controlling families use well-established ``central firms'' to acquire low-profitability entities through pyramidal structures; central firms trade at a discount to other group firms, with shareholders anticipating value-destroying acquisitions, though some underperformance may reflect selection rather than tunneling since low-profitability firms are placed in pyramids by design \citep{almeida_chaebol_formation}. The cash flow rights disparity between voting and economic ownership is substantially larger than previously reported, as studies using only publicly listed firms miss the amplifying role of non-public affiliates \citep{kim_chaebol_cf_rights}.

When a Chaebol-affiliated firm makes an acquisition, its own stock price falls on average while the controlling shareholder benefits through value transfer to other group firms \citep{bae_tunneling_value_added}. Business group firms also make systematically larger charitable donations, coordinated to tunnel resources from public affiliates (where family cash flow rights are lower) to private ones (where family stakes are higher) \citep{bae_business_groups_tunneling}.

The most causally credible evidence exploits Korea's 1999 reform imposing stricter board structure requirements on large firms (assets above 2 trillion won) relative to mid-sized firms: large firms whose controllers had incentives to tunnel earned strong positive abnormal returns around the reform, and panel regressions over 1998--2004 confirm that better governance moderates the negative effect of related-party transactions on firm value and raises the sensitivity of firm profitability to industry profitability \citep{black_governance_firm_value}. This natural experiment isolates the self-dealing channel and provides the identification that cross-sectional governance regressions cannot.

Concentrated ownership also distorts the broader governance environment. Among Chaebol-affiliated firms, only shareholder rights protection retains explanatory power for dividend payouts; board effectiveness, audit competency, and disclosure transparency lose independent influence once ownership structure is accounted for \citep{njoku_corporate_governance_dividends}. This pattern shows that standard improvements in board composition and disclosure are insufficient when concentrated ownership remains entrenched. The structural consequences include a persistent holding company discount \citep{suh_holding_company_discount}---unique to Korea and absent from Japan and the US---and the use of cross-shareholding arrangements as management entrenchment devices \citep{park_cross_shareholdings_takeover}.

\subsection{The Limits of Institutional Stewardship in Korea}
A first wave of policy responses sought to leverage the National Pension Service (NPS) as an activist institutional shareholder. NPS shareholding improves ESG performance and accounting and market outcomes among Korean manufacturing firms, with ESG improvements serving as the transmission channel \citep{kim_nps_esg_manufacturing}. NPS disclosures of ESG engagement intentions are associated with significant firm value gains in the following year, with effects largest for firms with the weakest internal governance---indicating that institutional activism substitutes for, rather than complements, internal governance mechanisms \citep{lee_esg_nps_firm_value}.

More direct forms of activism show limited reach. Markets do not react to NPS ``Vote No'' press announcements in the short run---a result inconsistent with developed-market shareholder activism findings---though firms that subsequently improve internal governance after a vote-no campaign do show higher valuation over time \citep{kim_nps_governance_firm_value}. The short-run market non-reaction underscores that NPS activism lacks enforcement credibility absent structural reform. Stewardship Code adoption by institutional blockholders is positively correlated with earnings quality, but multivariate regressions fail to establish statistical significance, attributed to a time lag before shareholder activism translates into measurable governance changes \citep{park_stewardship_code_earnings}.

A broader challenge emerges from a cross-country study covering 58 markets from 2001 to 2018: Korea's formal shareholder rights scores are not distinctively lower than peers, and markets with the largest shareholder return increases actually saw declines in average firm value \citep{kim_korea_discount_causes}. What distinguishes Korea is the low value relevance of fundamentals---profitability and growth are incorporated into prices to a far lesser degree than in comparable markets (51st of 58 countries)---suggesting a short-term investment culture prevents long-term growth expectations from being capitalized.

Together, institutional activism within the current legal framework is insufficient to close the P/B gap absent structural reform to shareholder enforcement.

\subsection{Top-Down Policy: Japan's Reforms as the Benchmark}
Contemporary literature has converged on top-down institutional reform as the necessary complement to institutional activism, drawing heavily on Japan's staged governance transformation. Korea's 2024 ``Corporate Value-Up'' program's voluntary structure produced low firm participation and a weak market response---the KOSPI closed weaker in 2024, the opposite of Japan's experience following mandatory TSE P/B directives \citep{lee_value_up_effectiveness}. Japan's reforms proceeded through a staged progression: the 2014 Stewardship Code, the 2015 Corporate Governance Code, and ultimately the 2023 TSE exchange-level mandate requiring firms with P/B below 1.0 to disclose capital cost improvement plans---each stage ratcheting enforcement rather than relying on voluntary compliance.

The 2025 Korean Commercial Act amendment---which expanded directors' fiduciary duties from the corporation to shareholders broadly---provides direct evidence that legal reform shifts valuation expectations: firms with the highest book-to-market ratios earned significantly positive abnormal returns at both the preliminary expectation date and the parliamentary passage date, with the differential reaction larger around formal passage than around the preceding presidential election, and no significant reaction among high controlling-shareholder-ownership firms \citep{kang_commercial_act_market_reactions}. These patterns suggest markets expect the amendment to have selective rather than uniform effects---larger where governance reform has the most room to reduce expropriation. Foreign investors act as asymmetric disciplinary monitors, correcting both over-investment in over-growing firms and underperformance in under-growing ones, implying that reforms improving market accessibility would amplify market discipline through the foreign investor channel \citep{joo_foreign_investors_value_up}.

The emerging consensus is that resolving the Korea Discount requires a transition from voluntary guidelines to comprehensive structural mandates. Korea's current trajectory mirrors Japan's pre-2023 soft-law phase. Our empirical analysis corroborates this view. Korea's relative P/B penalty against Japan deepened precisely during the period in which Japan escalated to exchange-level enforcement while Korea remained on a voluntary program. Critically, the same pattern holds when Japan is replaced by MSCI EM Asia as the comparator---ruling out Japan-specific post-reform appreciation as the mechanical driver---consistent with the view that hard structural mandates, not soft codes, are the critical threshold for sustained valuation recovery.
